\subsubsection{PackNet}

\gls{packnet} \cite{ll_packnet_2018} is a parameter-isolation \gls{cl} method
that allocates disjoint subsets of network weights to successive tasks. Rather
than penalising weight drift or replaying past data, \gls{packnet} physically
partitions parameters through iterative pruning, ensuring that no task can
overwrite the weights locked for a previous task.

\paragraph{Parameter Ownership}

Each prunable weight $\theta_i$ carries an owner label $o_i \in \{-1,\,0,\,1,
\ldots,T\}$, where $o_i = -1$ denotes a \emph{free} (unallocated) weight and
$o_i = t \geq 0$ denotes a weight permanently assigned to task $t$. A parallel
\emph{frozen buffer} $\hat\theta_i$ caches the value of every locked weight and
is reinstated after every optimiser step, preventing any subsequent training from
corrupting it. Non-prunable parameters---\gls{bn} affines and the final
classification head---are excluded from the ownership mechanism and remain freely
trainable throughout.

\paragraph{Free-Parameter Update}

While learning task $t$, only free parameters are updated. Gradients of all
locked weights are zeroed before the optimiser step,
%
\begin{equation}\label{eqn:packnet_grad_mask}
    g_i \leftarrow g_i \cdot \mathbf{1}[o_i < 0],
\end{equation}
%
and afterwards the frozen buffer is reinstated,
%
\begin{equation}\label{eqn:packnet_restore}
    \theta_i \leftarrow \hat\theta_i
    \quad \forall\, i : o_i \geq 0 ,
\end{equation}
%
so that locked weights cannot be perturbed by either the gradient step or any
momentum-based optimiser state.

\paragraph{Packing: Magnitude-Based Pruning}

At the end of task $t$, a fraction $\rho$ of the currently free weights
are promoted to permanent ownership and the remainder are returned to the free
pool. Given the free-weight index set $\mathcal{F} = \{i : o_i < 0\}$, the
$\lceil\rho\,|\mathcal{F}|\rceil$ weights of largest magnitude are selected,
%
\begin{equation}\label{eqn:packnet_topk}
    \mathcal{K}_t
    =
    \operatorname{top-}\!\lceil\rho\,|\mathcal{F}|\rceil_{i\in\mathcal{F}}
    \;|\theta_i| ,
\end{equation}
%
and locked by writing
%
\begin{equation}\label{eqn:packnet_lock}
    o_i \leftarrow t,
    \quad
    \hat\theta_i \leftarrow \theta_i
    \quad \forall\, i \in \mathcal{K}_t .
\end{equation}
%
Weights not selected ($\mathcal{F}\setminus\mathcal{K}_t$) are set to zero and
remain free for future tasks. The keep ratio $\rho = 1 - \texttt{prune\_perc}$ is
a hyperparameter. Because the free pool shrinks geometrically---task $t$ starts
with $|\mathcal{F}_t| = (1-\rho)^{t-1}|\theta|$ free weights---a fixed $\rho$
allocates the same \emph{fraction} of whatever capacity remains rather than the
same absolute count; earlier tasks therefore receive more parameters than later
ones.

\paragraph{Task-Aware Forward Pass}

During the forward pass for task $t$, only weights whose owner satisfies
$o_i \leq t$ are active; all later-task and still-free weights are temporarily
zeroed:
%
\begin{equation}\label{eqn:packnet_fwd}
    \tilde\theta_i = \theta_i \cdot \mathbf{1}[o_i \leq t] .
\end{equation}
%
The original values are restored immediately after the forward pass, so the stored
parameter tensor is not mutated. This makes the computation equivalent to applying
a binary task mask to every prunable tensor.

\paragraph{BatchNorm Statistics}

To prevent cross-task interference through shared normalisation layers, the
running statistics (mean, variance, and tracked-batch count) and affine parameters
of every \gls{bn} layer are snapshotted at the end of each task's training and
restored from the corresponding snapshot at the start of each subsequent training
episode or evaluation pass.

\paragraph{Inference}

\gls{packnet} uses a standard linear classifier at inference time. The task-aware
mask \eqref{eqn:packnet_fwd} is applied to activate only the parameters relevant
to the queried task. In \gls{til}, task identity is used to select the relevant
task logit slice. In \gls{cil}, prediction is made from the full logit vector
across all seen classes.
